%% chapters/assignment/ch3_ue.tex
%% 第3章　利用者均衡配分（UE）

% 第5節 活動ベース・スケジューリング・離散連続モデル

% ここで「経路選択だけではない」ことを示す。

% 5.1 活動ベースアプローチ
% 旅行は活動派生需要
% 1トリップずつではなく一日の活動スケジュールとしてみる
% 5.2 ABM
% 個人単位のルールベース / 効用ベースのシミュレーション
% 多主体・相互作用・制約表現に強い
% ただし推定可能性・識別可能性は弱くなりやすい
% 5.3 離散連続モデル
% MDCEV
% Tobit
% 消費量や滞在時間の内生化
% corner solution

% この節は、行動モデル章を「ただの離散選択モデル章」にしないために必要です。


\section{活動経路の記述：アクティビティモデル}\label{sec:abms}
本節では，活動経路の選択行動を記述する行動モデルについて解説する．
一般に活動経路を記述するモデルはアクティビティモデルと呼ばれる．